\chapter{Limites de fonctions et croissances comparées}
\section*{Exercice 12.1 -- Limites latérales}
À droite, le dénominateur tend vers zéro par valeurs positives : limite $+\infty$. À gauche, il est négatif : limite $-\infty$. La droite $x=2$ est asymptote verticale ; il n'y a pas de limite bilatérale unique.

\section*{Exercice 12.2 -- Signe à l’infini}
L'expression vaut $(|x|/x)\sqrt{1+4/x^2}$. La limite est $1$ à $+\infty$ et $-1$ à $-\infty$. Remplacer $\sqrt{x^2}$ par $x$ aurait perdu le second signe.

\section*{Exercice 12.3 -- Exponentielle dominante}
Poser $t=2x$ : $x^5e^{-2x}=t^5/(32e^t)\to0$ par croissance comparée. L'exposant positif $2x$ au dénominateur domine la puissance.

\section*{Exercice 12.4 -- Asymptote oblique}
La division donne $x+1+3/(x-1)$. Donc $y=x+1$ est asymptote aux deux infinis et $x=1$ verticale. La courbe est au-dessus de l'oblique pour $x>1$, au-dessous pour $x<1$.

\section*{Exercice 12.5 -- Limite et valeur}
Pour $x\ne1$, $f(x)=x+1$, donc la limite est $2$. Elle diffère de $f(1)=0$ : la fonction n'est pas continue. Changer la valeur en $1$ à $2$ donnerait le prolongement continu.

\section*{Exercice 12.6 -- Une oscillation écrasée}
Pour $x\ne0$, $|x\sin(1/x)|\le|x|\to0$. L'encadrement donne la limite, même si le facteur $\sin(1/x)$ n'en possède pas.
